% =====================================================================
% Colored Observation Noise with Randomized Spectral Exponent
% for Sim-to-Real Domain Randomization (Overleaf-compatible)
%
% Required Packages:
%   \usepackage{amsmath}
%   \usepackage{amssymb}
%   \usepackage{bm}
%   \usepackage{booktabs}
% =====================================================================

\subsection{Colored Observation Noise with Randomized Spectral Exponent}
\label{sec:colored_obs_noise}

\subsubsection{Problem Statement}

Domain randomization corrupts simulated observations with noise to bridge the Sim-to-Real gap.
A common approach adds independent and identically distributed (i.i.d.)\ white noise
$\boldsymbol{\epsilon}_t \sim \mathcal{N}(\boldsymbol{0}, \sigma^2 \boldsymbol{I})$
at each time step $t$.
However, real sensor errors exhibit temporal correlations:
encoder backlash, thermal drift, and communication jitter each produce noise
with a characteristic frequency spectrum.
An alternative is to use \emph{pink noise} ($1/f$ noise, spectral exponent $\beta=1$)~\cite{eberhard2023pink},
which introduces temporal correlation and better approximates real sensor behavior.

A remaining limitation is that the spectral exponent is fixed to $\beta = 1$.
In practice, different noise sources on the same robot produce different spectral shapes:
encoder quantization errors resemble $\beta \approx 0.5$,
servo backlash follows $\beta \approx 1.0$--$1.5$,
and thermal drift approaches $\beta \approx 2.0$.
A policy trained under a single fixed $\beta$ may overfit to that particular noise structure.

\subsubsection{Proposed Method: Per-Episode $\beta$ Randomization}

We propose randomizing the spectral exponent $\beta$ as a domain randomization parameter.
At each episode reset, every parallel environment independently samples
\begin{equation}
  \beta_i \sim \mathcal{U}(\beta_{\min},\, \beta_{\max})
  \label{eq:beta_sample}
\end{equation}
where $\beta_{\min}$ and $\beta_{\max}$ define the range of noise spectra encountered during training.
A colored noise sequence with power spectral density $S(f) \propto 1/f^{\beta_i}$ is then generated
using the Timmer--Koenig FFT algorithm~\cite{timmer1995}, with the key modification that the
amplitude scaling in the frequency domain uses the sampled exponent:
\begin{equation}
  s_k = f_k^{-\beta_i / 2}
  \label{eq:scaling}
\end{equation}
rather than the fixed $s_k = f_k^{-1/2}$ used in standard pink noise generation.

The generated unit-variance noise is applied selectively to sensor-derived observations
(joint positions and velocities), while command signals and previous actions remain uncorrupted:
\begin{equation}
  \tilde{\boldsymbol{o}}_t^{(\mathrm{sens})} = \boldsymbol{o}_t^{(\mathrm{sens})} + \sigma_{\mathrm{obs}} \, \boldsymbol{\eta}_t^{(\beta_i)}
  \label{eq:corruption}
\end{equation}
where $\boldsymbol{\eta}_t^{(\beta_i)}$ denotes a sample from the colored noise process
generated with exponent $\beta_i$.

\subsubsection{Training Scenarios}

To evaluate the effect of $\beta$ randomization, we define the following comparison scenarios
alongside the existing baselines:

\begin{table}[htbp]
  \centering
  \caption{Noise injection conditions for comparative evaluation}
  \label{tab:colored_scenarios}
  \begin{tabular}{lccc}
    \toprule
    Scenario & Action Noise & Obs.\ Noise & $\beta$ \\
    \midrule
    No-Noise (baseline) & -- & -- & -- \\
    Obs-Noise (pink) & -- & Pink & $1$ (fixed) \\
    \textbf{Obs-Colored-Noise (ours)} & -- & \textbf{Colored} & $\boldsymbol{\mathcal{U}(0.5, 1.5)}$ \\
    Both-Colored-Noise & Pink & Colored & action: $1$, obs: $\mathcal{U}(0.5, 1.5)$ \\
    \bottomrule
  \end{tabular}
\end{table}

\subsubsection{Hyperparameters}

\begin{table}[htbp]
  \centering
  \caption{Hyperparameters for colored observation noise}
  \label{tab:colored_hyperparams}
  \begin{tabular}{lll}
    \toprule
    Parameter & Symbol & Value \\
    \midrule
    Spectral exponent range & $[\beta_{\min},\, \beta_{\max}]$ & $[0.5,\, 1.5]$ \\
    Observation noise scale & $\sigma_{\mathrm{obs}}$ & $0.02$ \\
    Noise buffer length & $T$ & $1000$ steps \\
    \bottomrule
  \end{tabular}
\end{table}

The range $\beta \in [0.5, 1.5]$ spans from nearly-white noise ($\beta = 0.5$, weak correlation)
to strongly correlated noise ($\beta = 1.5$, approaching brown noise),
covering the diversity of real sensor characteristics.

% =====================================================================
% BibTeX References (add to your .bib file)
% =====================================================================
% @inproceedings{eberhard2023pink,
%   title     = {Pink Noise Is All You Need: Colored Noise Exploration in Deep Reinforcement Learning},
%   author    = {Eberhard, Onno and Hollenstein, Jakob and Pinneri, Cristina and Martius, Georg},
%   booktitle = {Proceedings of the Eleventh International Conference on Learning Representations (ICLR 2023)},
%   year      = {2023},
%   url       = {https://openreview.net/forum?id=hQ9V5QN27eS},
% }
% @article{timmer1995,
%   title   = {On generating power law noise},
%   author  = {Timmer, J. and Koenig, M.},
%   journal = {Astronomy and Astrophysics},
%   volume  = {300},
%   pages   = {707--710},
%   year    = {1995},
% }
